\documentclass[12pt, a4paper]{article}

%% ─── Packages ───────────────────────────────────────────────────────────────
\usepackage[a4paper, margin=2.5cm]{geometry}
\usepackage{amsmath, amssymb}
\usepackage{graphicx}
\usepackage{booktabs}
\usepackage{array}
\usepackage{caption}
\usepackage{subcaption}
\usepackage{hyperref}
\usepackage{fancyhdr}
\usepackage{titlesec}
\usepackage{xcolor}
\usepackage{enumitem}
\usepackage{siunitx}
\usepackage{float}
\usepackage{listings}
\usepackage{parskip}

\setlength{\headheight}{15pt}
\setlength{\parskip}{1em}

%% ─── Page Style ─────────────────────────────────────────────────────────────
\pagestyle{fancy}
\fancyhf{}
\rhead{EE568 -- Project 1}
\lhead{Elif Topaloğlu}
\cfoot{\thepage}
\renewcommand{\headrulewidth}{0.4pt}

%% ─── Hyperref Setup ─────────────────────────────────────────────────────────
\hypersetup{
    colorlinks=true,
    linkcolor=blue!60!black,
    urlcolor=blue!60!black,
    citecolor=blue!60!black,
}

%% ─── Section Formatting ─────────────────────────────────────────────────────
\titleformat{\section}{\large\bfseries}{\thesection\quad}{0em}{}
\titleformat{\subsection}{\normalsize\bfseries}{\thesubsection\quad}{0em}{}

%% ─── SI Units Setup ─────────────────────────────────────────────────────────
\sisetup{per-mode=symbol}

%% ─── MATLAB Listing Style ───────────────────────────────────────────────────
\lstdefinestyle{matlab}{
    language=Matlab,
    basicstyle=\ttfamily\small,
    keywordstyle=\color{blue}\bfseries,
    commentstyle=\color{green!50!black},
    stringstyle=\color{orange},
    numbers=left,
    numberstyle=\tiny\color{gray},
    numbersep=8pt,
    frame=single,
    breaklines=true,
    captionpos=b,
    tabsize=4,
}

\begin{document}

%%─── Title Page ──────────────────────────────────────────────────────────────
\begin{titlepage}
    \centering
    \begin{minipage}{0.25\textwidth}
        \centering
        
\includegraphics[height=3cm]{figures/eelogo.png} 
    \end{minipage}
    \hfill
    \begin{minipage}{0.25\textwidth}
        \centering
        
\includegraphics[height=3cm]{figures/Orta_Doğu_Teknik_Üniversitesi.png}
    \end{minipage}
    
    \vspace{2cm}
    {\LARGE \textbf{Middle East Technical University} \par}
    \vspace{1cm}
    {\large Department of Electrical and Electronics Engineering \par}
    \vspace{2cm}
    \rule{\linewidth}{0.5mm}\\[0.5cm]
    {\huge \bfseries EE568 -- Selected Topics in Electrical Machines \par}
    {\LARGE \bfseries Project 1: Torque in a Variable Reluctance Machine \par}
    \rule{\linewidth}{0.5mm}\\[1.5cm]
    
    \begin{tabular}{rl}
        \textbf{Author:}    & Elif Topaloğlu \\[0.5em]
        \textbf{Date:}      & \today \\[0.5em]
    \end{tabular}
    \vfill
    {\large \textbf{GitHub Repository:} \\[0.3em]
    \url{https://github.com/odtu/EE568/tree/master/Project1}}
\end{titlepage}
%%──────────────────────────────────────────────────────────






%%─── Table of Contents ───────────────────────────────────────────────────────
\tableofcontents
\newpage

%%─── 1. Introduction ─────────────────────────────────────────────────────────
\section{Introduction}

This report presents the analytical modeling and simulation analysis of a variable reluctance machine (VRM) as part of EE568 Project~1. Variable reluctance machines generate torque by minimizing the reluctance of the magnetic circuit; unlike permanent magnet or wound-rotor machines, they require no rotor excitation, making them robust and low-cost.

The main objectives of this project are:
\begin{enumerate}[leftmargin=*, label=\arabic*.]
    \item Derive analytical expressions for the reluctance, inductance, and torque as a function of rotor position.
    \item Model the system in a 2D finite element analysis (FEA) software using linear material properties and compare with the analytical results.
    \item Extend the FEA model to non-linear material properties and investigate the effects of saturation and fringing flux.
    \item Propose a control strategy to achieve continuous rotor rotation.
\end{enumerate}

The machine geometry for this project is defined on the course GitHub page. Some of the key parameters are summarized in Table~\ref{tab:params}.

\begin{table}[htbp]
    \centering
    \caption{Machine Parameters}
    \label{tab:params}
    \begin{tabular}{llc}
        \toprule
        \textbf{Parameter} & \textbf{Symbol} & \textbf{Value} \\
        \midrule
        Number of turns          & $N$         & 300 \\
        Coil current (DC)        & $I$         & \SI{2.5}{A} \\
        Core depth               & $d$         & \SI{25}{mm} \\
        Core width               & $w_c$       & \SI{15}{mm} \\
        Air-gap clearance        & $g$         & \SI{0.5}{mm} \\
        Coil winding area        & $A_{coil}$  & \SI{30}{mm} $\times$ \SI{10}{mm} \\
        \bottomrule
    \end{tabular}
\end{table}








\section{Q1 - Analytical Modelling}

\subsection{Assumptions}

The following simplifying assumptions are adopted throughout the analytical model:

\begin{enumerate}[leftmargin=*, label=\arabic*.]
    \item The core material is infinitely permeable ($\mu_r \to \infty$), so the core reluctance is negligible.
    \item Inductance varies sinusoidally with rotor position.
    \item Fringing flux at the air-gap is neglected.
    \item The magnetic system is linear (no saturation).
    \item Magnetic flux is uniformly distributed across the core cross-section.
\end{enumerate}







\subsection{Part (a) - Reluctance and Inductance Analysis}
By using the assumptions given, reluctance of the system can be calculated by the below formula:

\begin{equation}
    \mathcal{R} = \frac{l}{\mu  A_{core}}
    \label{eq:reluctance}
\end{equation}

where $l$ is the path length, $\mu$ is the permeability and $A_{core}$ is the core cross-sectional area. As the assumption states, core has infinitely permeable which makes the core reluctance zero. Thus, total reluctance will only consider the air gap reluctance. As the air-gap changes with respect to position, there will be a minimum and a maximum reluctance value. In below, all the parameters to find the reluctance and the maximum and minimum reluctance values is calculated.

The core cross-sectional area is:
\begin{equation}
    A_c = w_c \cdot d = 15 \times 10^{-3} \times 25 \times 10^{-3} = \SI{375e-6}{m^2}
\end{equation}

When the rotor tooth is fully aligned with the stator pole which is the minimum inductance position, the total air-gap path length consists of two single air-gaps:

\begin{equation}
    l_{g,\min} = 2g = 2 \times 0.5 = \SI{1}{mm}
\end{equation}

When the rotor is in the fully unaligned position, the flux path must also traverse the additional \SI{2}{mm} each side:
\begin{equation}
    l_{g,\max} = 2g + 2 \times \SI{2}{mm} = 1 + 4 = \SI{5}{mm}
\end{equation}

Substituting into Equation~\eqref{eq:reluctance}:

\begin{align}
    \mathcal{R}_{\min} &= \frac{1 \times 10^{-3}}{4\pi \times 10^{-7} \times 375 \times 10^{-6}}
                         = \SI{2.12e6}{H^{-1}} \\[6pt]
    \mathcal{R}_{\max} &= \frac{5 \times 10^{-3}}{4\pi \times 10^{-7} \times 375 \times 10^{-6}}
                         = \SI{1.06e7}{H^{-1}}
\end{align}

The self-inductance of the coil is related to the reluctance by:
\begin{equation}
    L = \frac{N^2}{\mathcal{R}}
    \label{eq:inductance}
\end{equation}

Therefore, by using Equation~\eqref{eq:inductance}:

\begin{align}
    L_{\max} &= \frac{300^2}{\mathcal{R}_{\min}} = \frac{90{,}000}{2.122 \times 10^6} \approx \SI{42.41}{mH} \\[6pt]
    L_{\min} &= \frac{300^2}{\mathcal{R}_{\max}} = \frac{90{,}000}{1.061 \times 10^7} \approx \SI{8.48}{mH}
\end{align}

\subsubsection{Inductance as a Function of Rotor Position}

The question further asks for the variation of the inductance with respect to the position under the assumption of sinusoidal variation. Since the machine has two poles due to variable reluctance, one mechanical revolution will produce two electrical cycles. Assuming a sinusoidal variation of inductance with the electrical angle $\theta$:

\begin{equation}
    L(\theta) = L_\text{avg} + L_\Delta \cos(2\theta)
    \label{eq:L_theta}
\end{equation}
where:
\begin{align}
    L_\text{avg} &= \frac{L_{\max} + L_{\min}}{2} = \frac{42.41 + 8.48}{2} \approx \SI{25.45}{mH} \\[4pt]
    L_\Delta     &= \frac{L_{\max} - L_{\min}}{2} = \frac{42.41 - 8.48}{2} \approx \SI{16.96}{mH}
\end{align}

The resulting inductance profile is plotted in Figure~\ref{fig:inductance}.


 \begin{figure}[H]
     \centering
     \includegraphics[width=0.82\textwidth]{figures/inductance_variation.eps}
     \caption{Inductance variation as a function of rotor position.}
     \label{fig:inductance}
 \end{figure}

The inductance reaches its peak of \SI{42.41}{mH} at the aligned positions (\ang{0} and \ang{180}) and its minimum of \SI{8.48}{mH} at the unaligned positions (\ang{90} and \ang{270}). 

\subsubsection{Numerical Results Summary}
Finally, numerical results for this part is summarized below.

\begin{table}[H]
    \centering
    \caption{Analytical Reluctance and Inductance Results}
    \label{tab:results_a}
    \begin{tabular}{lcc}
        \toprule
        \textbf{Quantity} & \textbf{Symbol} & \textbf{Value} \\
        \midrule
        Minimum reluctance  & $\mathcal{R}_{\min}$  & \SI{2.122e6}{H^{-1}} \\
        Maximum reluctance  & $\mathcal{R}_{\max}$  & \SI{1.061e7}{H^{-1}} \\
        Maximum inductance  & $L_{\max}$            & \SI{42.41}{mH} \\
        Minimum inductance  & $L_{\min}$            & \SI{8.48}{mH} \\
        \bottomrule
    \end{tabular}
\end{table}











\subsection{Part (b) -Torque under Constant DC Excitation}

\subsubsection{Torque Derivation}

To derive the torque expression, one should make use of the stored magnetic energy on an inductor. The stored magnetic energy can be expressed in terms of flux linkage as:

\begin{equation}
W = \int_{0}^{\lambda} i(\lambda)\, d\lambda
\end{equation}

For a linear magnetic system, the relationship between flux linkage and current is:

\begin{equation}
\lambda = L(\theta)\, i
\end{equation}

Thus, the current can be written as:

\begin{equation}
i = \frac{\lambda}{L(\theta)}
\end{equation}

Substituting into the energy expression:

\begin{equation}
W = \int_{0}^{\lambda} \frac{\lambda}{L(\theta)} \, d\lambda
\end{equation}

Since $L(\theta)$ is independent of $\lambda$ (linear system), it can be treated as constant:

\begin{equation}
W = \frac{1}{L(\theta)} \int_{0}^{\lambda} \lambda \, d\lambda
\end{equation}

\begin{equation}
W = \frac{1}{L(\theta)} \cdot \frac{1}{2} \lambda^2
\end{equation}

Substituting $\lambda = L(\theta)i$:

\begin{equation}
W = \frac{1}{2} L(\theta) i^2
\label{eq:energy}
\end{equation}

The electromagnetic torque is obtained from:

\begin{equation}
T = \frac{dW}{d\theta}
\end{equation}

\begin{equation}
T = \frac{d}{d\theta} \left( \frac{1}{2} L(\theta) i^2 \right)
\end{equation}

For constant current excitation:

\begin{equation}
T(\theta) = \frac{1}{2} i^2 \frac{dL(\theta)}{d\theta}
\label{eq:torque}
\end{equation}

Inductance variation is defined in the previous question. By taking the derivative, one can obtain the inductance variation as follows:

\begin{equation}
L(\theta) = L_{\text{avg}} + L_{\text{amp}} \cos(2\theta)
\end{equation}

\begin{equation}
\frac{dL}{d\theta} = -2 L_{\text{amp}} \sin(2\theta)
\label{eq:inductance_variation}
\end{equation}

By inserting the inductance variation into Equation~\eqref{eq:torque}, the final torque equation becomes:

\begin{equation}
T(\theta) = -i^2 L_{\text{amp}} \sin(2\theta)
\end{equation}

\subsubsection{Torque as a Function of Rotor Position}

The torque profile is plotted in Figure~\ref{fig:torque_variation} by using the derived equation.

\begin{figure}[H]
    \centering
    \includegraphics[width=0.82\textwidth]{figures/torque_variation.eps}
    \caption{Electromagnetic torque as a function of rotor position under constant DC excitation ($I = \SI{2.5}{A}$).}
    \label{fig:torque_variation}
\end{figure}

\subsubsection{Results Summary}

\begin{table}[htbp]
    \centering
    \caption{Torque Results (Constant DC Excitation)}
    \label{tab:results_b}
    \begin{tabular}{lcc}
        \toprule
        \textbf{Quantity} & \textbf{Symbol} & \textbf{Value} \\
        \midrule
        Maximum torque  & $T_{\max}$  & \SI{+0.11}{Nm} \\
        Minimum torque  & $T_{\min}$  & \SI{-0.11}{Nm} \\
        \bottomrule
    \end{tabular}
\end{table}

The torque is sinusoidal at twice the rotor frequency. It is positive in the regions where inductance is increasing (rotor moving towards alignment) and negative where inductance is decreasing. Under constant DC excitation, the average torque over a full revolution is zero, since the positive and negative half-cycles cancel.









\subsection{Part (c) - Non-Linear Effects and Model Improvements}

\subsubsection{Fringing Flux}

In the simple analytical model, flux is assumed to flow only through the face of the stator pole directly across the air-gap, ignoring the flux that flows outward around the pole edges. In practice, this fringing flux can significantly increase the effective cross-sectional area, particularly when the air-gap is large. One common approach to consider the fringing effects is to add the gap length to each linear dimension of the pole face to obtain an effective area:

\begin{equation}
    A_\text{eff} = (w_c + g)(d + g)
\end{equation}

This correction, as detailed in \cite{hurley2013}, is applied separately for each rotor position and reduces the reluctance, raising the predicted inductance closer to FEA results.

\subsubsection{Non-Uniform Flux Distribution}

The assumption of uniform flux distribution across the core is valid when the core is far from saturation conditions or the dimensions are such that the flux path is relatively uniform. A more accurate treatment to include would divide the core and air-gap regions into parallel flux tubes and solve for the flux in each tube independently, accounting for the variation of permeability in a saturating core. This would require numerical integration across the core cross-section and is more complex, but it would yield a more accurate prediction of inductance and torque, especially near saturation.

\subsubsection{Saturation and Non-Linear Permeability}

When the flux density in the core approaches the saturation flux density ($B_\text{sat} \approx \SI{1.5}{T}$ for typical electrical steel), the incremental permeability drops sharply. In this regime, the inductance is no longer constant for a given position. It becomes a function of both $\theta$ and $I$. For this condition, the co-energy must be computed numerically from the magnetisation curve:

\begin{equation}
    W'(I,\theta) = \int_0^I \lambda(i, \theta)\, di
\end{equation}

and the torque obtained as $T = \partial W' / \partial \theta \big|_{I=\text{const}}$.

To calculate this effect, the B-H curve of the selected lamination material to be included in the analysis.










\section{Q2 - 2D FEA Modelling (Linear Materials)}

In this section, the variable reluctance machine will be modelled in a 2D finite element analysis (FEA) software using linear material properties. The results will be compared with the analytical model derived in Q1 to validate the assumptions and identify any discrepancies. Main outputs from the FEA simulations will include flux density vectors, inductance values at different rotor positions, stored magnetic energy, and torque profiles. MAchine is generated in ANSYS Maxwell software by using the geometry provided on the course GitHub page. As the core material, M-19 electrical steel is selected due to its typical use in such applications and low eddy current losses. Material properties of the M-19 steel is available on ANSYS and the properties are presented in the appendix. In the linear part of the B-H curve, material has the relative permeability of $\mu_r \approx 5000$ which is used in the linear FEA simulations. Results for the analysis will be presented in the following subsections. Simulation model is presented in Figure~\ref{fig:fea_model}.

\begin{figure}[H]
    \centering
    \includegraphics[width=0.7\textwidth]{figures/fea_model.png}
    \caption{2D FEA model of the variable reluctance machine in ANSYS Maxwell.}
    \label{fig:fea_model}
\end{figure}

\subsection{Part (a) - Flux Density Vectors at Different Rotor Positions}

In this part, flux density vectors are plotted for three different rotor positions: aligned (0°), intermediate (45°), and unaligned (90°). The flux density distribution will illustrate how the magnetic field lines concentrate in the aligned position and spread out in the unaligned position, confirming the assumptions made in the analytical model. The flux density plots will be presented in Figures \ref{fig:flux_0deg}, \ref{fig:flux_45deg}, and \ref{fig:flux_90deg}.

As can be seen from the flux density plots, the flux lines are highly concentrated at the starting position (0°) where the rotor tooth is fully aligned with the stator pole, resulting in a high flux density. At the intermediate position (45°), the flux lines begin to spread out, indicating a decrease in flux density. 

Finally, at the unaligned position (90°), the flux lines are widely spread, confirming the significant reduction in flux density and validating the assumptions of the analytical model. As the air gap increases with the rotor position, the reluctance increases which decreases the flux density while increasing the fringing fluxes across the air gap. These fringing fields can be observed in the flux density plots.

\begin{figure}[H]
    \centering
    \includegraphics[width=0.7\textwidth]{figures/flux_0deg.png}
    \caption{Flux density at 0° for linear material}
    \label{fig:flux_0deg}
\end{figure}

\hfill
\begin{figure}[H]
    \centering
    \includegraphics[width=0.7\textwidth]{figures/flux_45deg.png}
    \caption{Flux density at 45° for linear material}
    \label{fig:flux_45deg}
\end{figure}
\hfill
\begin{figure}[H]
    \centering
    \includegraphics[width=0.7\textwidth]{figures/flux_90deg.png}
    \caption{Flux density at 90° for linear material}
    \label{fig:flux_90deg}
\end{figure}




\subsection{Part (b) - Inductance and Stored Energy}

For the same three rotor positions, the inductance values and stored magnetic energy will be calculated from the FEA results. The inductance can be obtained from ANSYS itself by assigning the appropriate inductance parameters as the matrix, while the stored energy can be calculated using the co-energy method using Equation~\eqref{eq:energy}. These values will be compared with the analytical predictions to assess the accuracy of the analytical calculations and identify any discrepancies due to the assumptions made. Inductance results will be tabulated in Table~\ref{tab:inductance_fea} and the stored energy values will be presented in Table~\ref{tab:energy_fea}. 

When the results are compared with the analytical model, it is expected that the FEA results will show higher inductance values at the aligned position due to the fringing flux which increases the effective area. In FEA, inductance is calculated as 49.4 mH at 0° which is higher than the analytical value of 42.41 mH. At the intermediate position (45°), the inductance is calculated as 32.2 mH in FEA which is also higher than the analytical value of 25.45 mH. Finally, at the unaligned position (90°), the inductance from FEA is 16.5 mH which is again higher than the analytical value of 8.48 mH. This discrepancy can be attributed to the fringing flux effects which are not accounted for in the analytical model, confirming that fringing flux significantly increases the effective area and thus the inductance, especially at larger air gaps which is the reson for the larger discrepancy at the unaligned position. 

Additionally, as the analytical model assumes a uniform flux distribution, it underestimates the inductance at positions where the flux is more concentrated, while FEA captures the non-uniform distribution and thus predicts higher inductance values. Flux density plots from the previous section also support this conclusion. 

\begin{table}[htbp]
    \centering
    \caption{Inductance Values from FEA for Linear Material}
    \label{tab:inductance_fea}
    \begin{tabular}{|c|c|c|}
        \hline
        Rotor Position & Inductance (mH) \\
        \hline
        0° & 49.4 \\
        45° & 32.2 \\
        90° & 16.5 \\
        \hline
    \end{tabular}
\end{table}

\begin{table}[htbp]
    \centering
    \caption{Stored Energy Values from FEA for Linear Material}
    \label{tab:energy_fea}
    \begin{tabular}{|c|c|c|}
        \hline
        Rotor Position & Stored Energy (J) \\
        \hline
        0° & 0.154 \\
        45° & 0.101 \\
        90° & 0.052 \\
        \hline
    \end{tabular}
\end{table}




\subsection{Part (c) - Torque Variation with respect to Rotor Angle and Comparison with Analytical Model}

In FEA, rotor is rotated in small increments of the angle (e.g., every 1°) and the torque is obtained at each position. The resulting torque profile is plotted as a function of rotor angle and presented in Figure~\ref{fig:torque_variation_fea}.  The torque profile was assumed to have a sinusoidal variation in the analytical model. However, the FEA results may show deviations from this idealized shape as the rotor has different geometries at different angles that can cause non-sinusoidal variations in the inductance and thus the torque. Thus, FEA results show that the sinusoidal waveform assumption is not entirely accurate, especially near the aligned and unaligned positions where the flux distribution changes more abruptly. 

\begin{figure}[h]
    \centering
    \includegraphics[width=0.82\textwidth]{figures/Torque vs Rotation Angle Plot.png}
    \caption{Torque variation with respect to rotor angle from FEA simulations for linear material.}
    \label{fig:torque_variation_fea}
\end{figure}

Also, torque waveform is not symmetric as the analytical model assumes. This is because the analytical model assumes a uniform flux distribution and neglects fringing effects, which makes the change of inductance with respect to rotor angle symmetric. However, in FEA, the flux distribution is non-uniform and fringing effects are included, which leads to an asymmetric torque profile. 

Additionally, peak torque from the FEA is approximetely \SI{81.8}{mNm} which is lower than the analytical value of \SI{110}{mNm}. The reason is again the fringing flux which increases the inductance for the unaligned position more than the aligned position, thus reducing the maximum torque values.



\section{Q3 - 2D FEA Modelling (Non-Linear Materials)}

In this section, the FEA model will be extended to include non-linear material properties, specifically the B-H curve of the M-19 electrical steel. The effects of saturation and fringing flux will be investigated by comparing the results with the linear FEA model and the analytical predictions. The main outputs will include flux density vectors, inductance values, stored energy, and torque profiles at different rotor positions. The results will be analyzed to understand how non-linearities affect the machine performance and to identify any limitations of the linear assumptions made in the analytical model.


\subsection{Part (a) - Flux Density Vectors at Different Rotor Positions}

As before, flux density vectors are plotted for three different rotor positions: aligned (0°), intermediate (45°), and unaligned (90°) for the non-linear material. The flux density plots will be presented in Figures \ref{fig:flux_0deg_NL}, \ref{fig:flux_45deg_NL}, and \ref{fig:flux_90deg_NL}.

\begin{figure}[H]
    \centering
    \includegraphics[width=0.7\textwidth]{figures/flux_0deg_NL.png}
    \caption{Flux density at 0° for non-linear material}
    \label{fig:flux_0deg_NL}
\end{figure}

\hfill
\begin{figure}[H]
    \centering
    \includegraphics[width=0.7\textwidth]{figures/flux_45deg_NL.png}
    \caption{Flux density at 45° for non-linear material}
    \label{fig:flux_45deg_NL}
\end{figure}
\hfill
\begin{figure}[H]
    \centering
    \includegraphics[width=0.7\textwidth]{figures/flux_90deg_NL.png}
    \caption{Flux density at 90° for non-linear material}
    \label{fig:flux_90deg_NL}
\end{figure}

When the figures are compared with the linear case, there is no significant difference in the flux density distribution for all the positions since the core is not saturated and the flux density is below the saturation level. If the current is increased further, the flux density would approach the saturation level and increase the non-linearity effects, which would lead to a more significant difference in the flux density distribution between the linear and non-linear cases. But for this current value and the steel material, the flux density is still in the linear region of the B-H curve, thus the flux density distribution is similar to the linear case. 


\subsection{Part (b) - Inductance and Stored Energy}

For the inductance and the stored energy, the results are similar to the linear case since the core is not saturated. The inductance values at different rotor positions are tabulated in Table~\ref{tab:inductance_fea_NL} and the stored energy values are presented in Table~\ref{tab:energy_fea_NL}.

\begin{table}[htbp]
    \centering
    \caption{Inductance Values from FEA for Non-Linear Material}
    \label{tab:inductance_fea_NL}
    \begin{tabular}{|c|c|c|}
        \hline
        Rotor Position & Inductance (mH) \\
        \hline
        0° &  \\
        45° &  \\
        90° &  \\
        \hline
    \end{tabular}
\end{table}

\begin{table}[htbp]
    \centering
    \caption{Stored Energy Values from FEA for Non-Linear Material}
    \label{tab:energy_fea_NL}
    \begin{tabular}{|c|c|c|}
        \hline
        Rotor Position & Stored Energy (J) \\
        \hline
        0° & 0. \\
        45° & 0. \\
        90° & 0. \\
        \hline
    \end{tabular}
\end{table}


\subsection{Part (c) - Torque Variation with respect to Rotor Angle and Comparison with Analytical Model}

Finally, the torque variation with respect to rotor angle is plotted for the non-linear material and compared with the analytical model. The torque profile is presented in Figure~\ref{fig:torque_variation_fea_NL}. The torque profile is similar to the linear case since the core is not saturated. However, if the current is increased further, the torque profile would deviate more significantly from the linear case due to the saturation effects, which would lead to a reduction in the maximum torque and a distortion of the torque waveform. But for this current value and the steel material, the torque profile is similar to the linear case.


\begin{figure}[htbp]
    \centering
    \includegraphics[width=0.7\textwidth]{figures/Torque vs Rotation Angle Plot.png}
    \caption{Torque variation with respect to rotor angle from FEA simulations for Non-Linear Material.}
    \label{fig:torque_variation_fea_NL}
\end{figure}


\subsection{Part (d) - Comparison: Linear vs. Non-Linear, Fringing and Saturation Effects}










\section{Q4 - Control Method for Continuous Rotation}

\textit{This section will be completed.}












\section{Conclusion}

The analytical model of the variable reluctance machine was developed under linearised assumptions. The maximum and minimum inductances were found to be \SI{42.41}{mH} and \SI{8.48}{mH}, corresponding to the aligned and unaligned rotor positions respectively. The peak torque under \SI{2.5}{A} DC excitation is $\pm\SI{0.106}{Nm}$, varying sinusoidally at twice the rotor frequency.

FEA simulations (Q2, Q3) and the control strategy (Q4) will be added as the project progresses.










\appendix
\section{Material Properties}

This appendix will list the B-H curve data and material properties of the electrical steel lamination selected for the FEA simulations in Q2 and Q3.

\section{MATLAB Source Code}

The MATLAB script used for the analytical calculations and plots in Q1 is provided below.

%\lstinputlisting[style=matlab, caption={MATLAB script for Q1 Analytical Modelling}]{analytical_model.m}

%%─── References ──────────────────────────────────────────────────────────────
\begin{thebibliography}{9}
\bibitem{hurley2013} 
W. G. Hurley and W. H. Wölfle, 
\textit{Transformers and Inductors for Power Electronics: Theory, Design and Applications}, 
1st ed. Chichester, UK: John Wiley \& Sons, 2013. 
ISBN: 9781118544679.
\end{thebibliography}

\end{document}